\documentclass[12pt]{article}
\usepackage{amsmath, amssymb}
\usepackage[a4paper,margin=2.5cm]{geometry}

\begin{document}
	
	\title{Study Notes: Determinants}
	\author{Antony Lotter}
	\date{\today}
	\maketitle
	
	\section{Determinants}
	
	\subsection{Definition and Notation}
	
	Let $A$ be an $n \times n$ (square) matrix with real entries.  
	The \emph{determinant} of $A$ is a scalar, denoted $\det(A)$ or $|A|$, that encodes important algebraic and geometric information about the linear transformation represented by $A$.
	
	Key viewpoints:
	\begin{itemize}
		\item Algebraically, $\det(A) \neq 0$ if and only if $A$ is invertible.
		\item Geometrically, for a linear map $T_A:\mathbb{R}^n \to \mathbb{R}^n$ given by $T_A(\mathbf{x}) = A\mathbf{x}$, the absolute value $|\det(A)|$ is the factor by which $A$ scales $n$-dimensional volume.
		\item The sign of $\det(A)$ indicates whether orientation is preserved ($\det(A) > 0$) or reversed ($\det(A) < 0$).
	\end{itemize}
	
	\subsection{Small-Size Determinants}
	
	\paragraph{2$\times$2 case.}
	For
	\[
	A = \begin{pmatrix}
		a & b \\
		c & d
	\end{pmatrix},
	\]
	the determinant is
	\[
	\det(A) = ad - bc.
	\]
	
	\paragraph{3$\times$3 case via cofactor expansion.}
	For
	\[
	A = \begin{pmatrix}
		a_{11} & a_{12} & a_{13} \\
		a_{21} & a_{22} & a_{23} \\
		a_{31} & a_{32} & a_{33}
	\end{pmatrix},
	\]
	a common formula (expansion along the first row) is
	\[
	\det(A)
	=
	a_{11}
	\begin{vmatrix}
		a_{22} & a_{23} \\
		a_{32} & a_{33}
	\end{vmatrix}
	-
	a_{12}
	\begin{vmatrix}
		a_{21} & a_{23} \\
		a_{31} & a_{33}
	\end{vmatrix}
	+
	a_{13}
	\begin{vmatrix}
		a_{21} & a_{22} \\
		a_{31} & a_{32}
	\end{vmatrix},
	\]
	where each $2\times 2$ determinant is computed as in the $2\times 2$ case.
	
	\subsection{Minors, Cofactors, and Cofactor Expansion}
	
	Let $A = (a_{ij})$ be an $n\times n$ matrix.
	
	\begin{itemize}
		\item The \emph{minor} $M_{ij}$ of $a_{ij}$ is the determinant of the $(n-1)\times(n-1)$ matrix obtained by deleting row $i$ and column $j$ from $A$.
		\item The \emph{cofactor} $C_{ij}$ of $a_{ij}$ is
		\[
		C_{ij} = (-1)^{i+j} M_{ij}.
		\]
	\end{itemize}
	
	The determinant of $A$ can be computed by \emph{cofactor expansion} along any row $i$:
	\[
	\det(A) = \sum_{j=1}^{n} a_{ij} C_{ij},
	\]
	or along any column $j$:
	\[
	\det(A) = \sum_{i=1}^{n} a_{ij} C_{ij}.
	\]
	
	This formula is known as \emph{Laplace expansion}. A practical strategy is to expand along a row or column that contains many zeros to minimize computation.
	
	\subsection{Row Operations and Determinant}
	
	Elementary row operations affect the determinant in specific ways:
	
	\begin{enumerate}
		\item Swapping two rows multiplies the determinant by $-1$.
		\item Multiplying a row by a nonzero scalar $k$ multiplies the determinant by $k$.
		\item Adding a multiple of one row to another row does \emph{not} change the determinant.
	\end{enumerate}
	
	These properties allow efficient computation of determinants by row-reducing a matrix to an upper triangular form and tracking how the determinant changes.
	
	\bigskip
	
	% ===============================
	% Identity Matrices
	% ===============================
	
	\section{Identity Matrices}
	
	\subsection{Definition}
	
	The $n \times n$ \emph{identity matrix} $I_n$ is the matrix with ones on the main diagonal and zeros elsewhere:
	\[
	I_n = \begin{pmatrix}
		1 & 0 & \cdots & 0 \\
		0 & 1 & \cdots & 0 \\
		\vdots & \vdots & \ddots & \vdots \\
		0 & 0 & \cdots & 1
	\end{pmatrix}.
	\]
	
	\subsection{Key Properties}
	
	\begin{itemize}
		\item $I_n$ is the multiplicative identity for matrix multiplication:
		\[
		I_n A = A I_n = A
		\]
		for any $n \times n$ matrix $A$.
		\item The determinant of the identity matrix is $1$:
		\[
		\det(I_n) = 1.
		\]
		\item The identity matrix represents the linear transformation that leaves every vector unchanged.
	\end{itemize}
	
	\bigskip
	
	% ===============================
	% Inverse Matrices via Cofactors
	% ===============================
	
	\section{Inverse Matrices (Cofactor/Adjoint Method)}
	
	\subsection{Definition}
	
	An $n\times n$ matrix $A$ is called \emph{invertible} (or \emph{nonsingular}) if there exists a matrix $A^{-1}$ such that
	\[
	A A^{-1} = A^{-1} A = I_n.
	\]
	
	If no such matrix exists, $A$ is called \emph{singular}.
	
	\subsection{Adjugate (Adjoint) Matrix}
	
	The \emph{adjugate} (or \emph{adjoint}) of $A$, denoted $\operatorname{adj}(A)$, is defined as
	\[
	\operatorname{adj}(A) = C^{T},
	\]
	where $C$ is the matrix of cofactors:
	\[
	C = (C_{ij}), \qquad C_{ij} = (-1)^{i+j} M_{ij}.
	\]
	
	A fundamental identity relating $A$, its adjugate, and the identity matrix is
	\[
	A \cdot \operatorname{adj}(A) = \operatorname{adj}(A) \cdot A = \det(A) \cdot I_n.
	\]
	
	\subsection{Inverse via Adjugate}
	
	If $\det(A) \neq 0$, then $A$ is invertible and the inverse can be expressed using the adjugate:
	\[
	A^{-1} = \frac{1}{\det(A)} \operatorname{adj}(A).
	\]
	
	This formula is exact but becomes computationally expensive for large $n$ (because it requires many cofactors). For small matrices (e.g.\ $2\times 2$ or $3\times 3$) it is often reasonable.
	
	\subsection{Example (2$\times$2)}
	
	Let
	\[
	A = 
	\begin{pmatrix}
		1 & 2 \\
		3 & 4
	\end{pmatrix}.
	\]
	First compute the determinant:
	\[
	\det(A) = 1\cdot 4 - 2\cdot 3 = -2.
	\]
	
	Cofactors:
	\[
	C_{11} = 4,\quad
	C_{12} = -3,\quad
	C_{21} = -2,\quad
	C_{22} = 1.
	\]
	
	Cofactor matrix:
	\[
	C = 
	\begin{pmatrix}
		4 & -3 \\
		-2 & 1
	\end{pmatrix}.
	\]
	
	Adjugate:
	\[
	\operatorname{adj}(A) = C^T =
	\begin{pmatrix}
		4 & -2 \\
		-3 & 1
	\end{pmatrix}.
	\]
	
	Inverse:
	\[
	A^{-1} = \frac{1}{\det(A)} \operatorname{adj}(A)
	= \frac{1}{-2}
	\begin{pmatrix}
		4 & -2 \\
		-3 & 1
	\end{pmatrix}
	=
	\begin{pmatrix}
		-2 & 1 \\
		\frac{3}{2} & -\frac{1}{2}
	\end{pmatrix}.
	\]
	
	\bigskip
	
	% ===============================
	% Transpose Matrices
	% ===============================
	
	\section{Transpose Matrices}
	
	\subsection{Definition}
	
	Given an $m \times n$ matrix $A = (a_{ij})$, its \emph{transpose} $A^T$ is the $n \times m$ matrix obtained by interchanging rows and columns:
	\[
	(A^T)_{ij} = a_{ji}.
	\]
	
	In words, the $i$-th row of $A$ becomes the $i$-th column of $A^T$ (and vice versa).
	
	\subsection{Basic Properties}
	
	For matrices $A$ and $B$ of compatible sizes, and scalars $\alpha,\beta$:
	\begin{itemize}
		\item $(A^T)^T = A$.
		\item $(A + B)^T = A^T + B^T$.
		\item $(\alpha A)^T = \alpha A^T$.
		\item $(AB)^T = B^T A^T$.
	\end{itemize}
	
	If $A$ is $n\times n$, then so is $A^T$. In that case, the determinant satisfies
	\[
	\det(A^T) = \det(A).
	\]
	
	This means that transposing a square matrix does not change its determinant.
	
	\bigskip
	
	% ===============================
	% Singular and Unique Solutions
	% ===============================
	
	\section{Singular Matrices and Unique Solutions}
	
	\subsection{Singular vs Nonsingular}
	
	For an $n\times n$ matrix $A$:
	\begin{itemize}
		\item $A$ is \emph{nonsingular} (invertible) if $\det(A) \neq 0$.
		\item $A$ is \emph{singular} (non-invertible) if $\det(A) = 0$.
	\end{itemize}
	
	\subsection{Relationship to Solutions of $A\mathbf{x} = \mathbf{b}$}
	
	Consider the linear system
	\[
	A\mathbf{x} = \mathbf{b},
	\]
	where $A$ is $n\times n$.
	
	\begin{itemize}
		\item If $\det(A) \neq 0$ (nonsingular $A$), then:
		\begin{itemize}
			\item $A$ has an inverse $A^{-1}$.
			\item The system has a \emph{unique} solution:
			\[
			\mathbf{x} = A^{-1}\mathbf{b}.
			\]
		\end{itemize}
		\item If $\det(A) = 0$ (singular $A$), then:
		\begin{itemize}
			\item $A$ does not have an inverse.
			\item The system either has no solution or infinitely many solutions.
			\item There is \emph{never} a unique solution in this case.
		\end{itemize}
	\end{itemize}
	
	Thus, for square systems, the determinant acts as a test for the existence of a unique solution.
	
	\bigskip
	
	% ===============================
	% Cramer's Rule
	% ===============================
	
	\section{Cramer's Rule}
	
	\subsection{Statement}
	
	Cramer's rule provides a formula for the unique solution of a square linear system using determinants.
	
	Let $A$ be an $n\times n$ matrix with $\det(A)\neq 0$, and consider the system
	\[
	A\mathbf{x} = \mathbf{b},
	\quad
	\mathbf{x} =
	\begin{pmatrix}
		x_1 \\
		x_2 \\
		\vdots \\
		x_n
	\end{pmatrix},\quad
	\mathbf{b} =
	\begin{pmatrix}
		b_1 \\
		b_2 \\
		\vdots \\
		b_n
	\end{pmatrix}.
	\]
	
	For each $i \in \{1,\dots,n\}$, define the matrix $A_i$ by replacing the $i$-th column of $A$ with the vector $\mathbf{b}$:
	\[
	A_i = \begin{pmatrix}
		\text{col}_1(A) & \cdots & \text{col}_{i-1}(A) & \mathbf{b} & \text{col}_{i+1}(A) & \cdots & \text{col}_n(A)
	\end{pmatrix}.
	\]
	
	Then Cramer's rule states:
	\[
	x_i = \frac{\det(A_i)}{\det(A)}.
	\]
	
	\subsection{Conditions and Interpretation}
	
	\begin{itemize}
		\item The formula is valid only if $\det(A) \neq 0$ (so $A$ is invertible and the solution is unique).
		\item Each component $x_i$ is expressed as a ratio of determinants, showing directly how the right-hand side $\mathbf{b}$ affects each variable.
		\item Although conceptually elegant, Cramer's rule is inefficient for large systems because computing many determinants is expensive. It is mainly used for small systems or theoretical arguments.
	\end{itemize}
	
	\subsection{Example: 2$\times$2 System}
	
	Solve
	\[
	\begin{cases}
		x + 2y = 5, \\
		3x + 4y = 6.
	\end{cases}
	\]
	
	Coefficient matrix and vector:
	\[
	A = 
	\begin{pmatrix}
		1 & 2 \\
		3 & 4
	\end{pmatrix}, \quad
	\mathbf{b} =
	\begin{pmatrix}
		5 \\
		6
	\end{pmatrix}.
	\]
	
	Compute the determinant of $A$:
	\[
	\det(A) = 1\cdot4 - 2\cdot3 = -2 \neq 0,
	\]
	so there is a unique solution.
	
	\paragraph{Compute $x_1 = x$.}
	Replace the first column by $\mathbf{b}$:
	\[
	A_1 =
	\begin{pmatrix}
		5 & 2 \\
		6 & 4
	\end{pmatrix}.
	\]
	Then
	\[
	\det(A_1) = 5\cdot 4 - 2\cdot 6 = 20 - 12 = 8.
	\]
	So
	\[
	x = x_1 = \frac{\det(A_1)}{\det(A)} = \frac{8}{-2} = -4.
	\]
	
	\paragraph{Compute $x_2 = y$.}
	Replace the second column by $\mathbf{b}$:
	\[
	A_2 =
	\begin{pmatrix}
		1 & 5 \\
		3 & 6
	\end{pmatrix}.
	\]
	Then
	\[
	\det(A_2) = 1\cdot 6 - 5\cdot 3 = 6 - 15 = -9.
	\]
	So
	\[
	y = x_2 = \frac{\det(A_2)}{\det(A)} = \frac{-9}{-2} = \frac{9}{2}.
	\]
	
	Thus the unique solution is
	\[
	(x,y) = \left(-4,\;\frac{9}{2}\right).
	\]
	
	\section*{Determinant of a $2\times 2$ matrix}
	
	Consider
	\[
	A = \begin{bmatrix}
		a & b \\
		c & d
	\end{bmatrix}.
	\]
	
	The determinant of $A$ is
	\[
	\det(A) = \begin{vmatrix}
		a & b \\
		c & d
	\end{vmatrix}
	= ad - bc.
	\]
	
	\bigskip
	
	\section*{Determinant of a $3\times 3$ matrix}
	
	Let
	\[
	A = \begin{bmatrix}
		a_{11} & a_{12} & a_{13} \\
		a_{21} & a_{22} & a_{23} \\
		a_{31} & a_{32} & a_{33}
	\end{bmatrix}.
	\]
	
	A cofactor expansion along the first row gives
	\[
	\det(A) =
	a_{11}
	\begin{vmatrix}
		a_{22} & a_{23} \\
		a_{32} & a_{33}
	\end{vmatrix}
	-
	a_{12}
	\begin{vmatrix}
		a_{21} & a_{23} \\
		a_{31} & a_{33}
	\end{vmatrix}
	+
	a_{13}
	\begin{vmatrix}
		a_{21} & a_{22} \\
		a_{31} & a_{32}
	\end{vmatrix}.
	\]
	
	Using the $2\times 2$ formula inside, this becomes
	\[
	\det(A) =
	a_{11}(a_{22}a_{33} - a_{23}a_{32})
	-
	a_{12}(a_{21}a_{33} - a_{23}a_{31})
	+
	a_{13}(a_{21}a_{32} - a_{22}a_{31}).
	\]
	
	Equivalently (Sarrus' rule):
	\[
	\det(A) =
	a_{11}a_{22}a_{33}
	+
	a_{12}a_{23}a_{31}
	+
	a_{13}a_{21}a_{32}
	-
	a_{13}a_{22}a_{31}
	-
	a_{11}a_{23}a_{32}
	-
	a_{12}a_{21}a_{33}.
	\]
	
	\bigskip
	
	\section*{Determinant of a $4\times 4$ matrix}
	
	Let
	\[
	A =
	\begin{bmatrix}
		a_{11} & a_{12} & a_{13} & a_{14} \\
		a_{21} & a_{22} & a_{23} & a_{24} \\
		a_{31} & a_{32} & a_{33} & a_{34} \\
		a_{41} & a_{42} & a_{43} & a_{44}
	\end{bmatrix}.
	\]
	
	The cofactor $C_{1j}$ of the entry $a_{1j}$ is
	\[
	C_{1j} = (-1)^{1+j} M_{1j},
	\]
	where $M_{1j}$ is the determinant of the $3\times 3$ minor obtained by deleting row 1 and column $j$.
	
	A cofactor expansion along the first row is
	\[
	\det(A) =
	a_{11} C_{11}
	+
	a_{12} C_{12}
	+
	a_{13} C_{13}
	+
	a_{14} C_{14}.
	\]
	
	Explicitly,
	\[
	\det(A) =
	a_{11}
	\begin{vmatrix}
		a_{22} & a_{23} & a_{24} \\
		a_{32} & a_{33} & a_{34} \\
		a_{42} & a_{43} & a_{44}
	\end{vmatrix}
	-
	a_{12}
	\begin{vmatrix}
		a_{21} & a_{23} & a_{24} \\
		a_{31} & a_{33} & a_{34} \\
		a_{41} & a_{43} & a_{44}
	\end{vmatrix}
	+
	a_{13}
	\begin{vmatrix}
		a_{21} & a_{22} & a_{24} \\
		a_{31} & a_{32} & a_{34} \\
		a_{41} & a_{42} & a_{44}
	\end{vmatrix}
	-
	a_{14}
	\begin{vmatrix}
		a_{21} & a_{22} & a_{23} \\
		a_{31} & a_{32} & a_{33} \\
		a_{41} & a_{42} & a_{43}
	\end{vmatrix}.
	\]
	
	Each $3\times 3$ determinant can be expanded using the $3\times 3$ formulas above.
	
	\bigskip
	
	\newpage
	
	\section*{Useful properties}
	
	For any $n\times n$ matrix $A$:
	\begin{itemize}
		\item $\det(A) = 0$ if and only if $A$ is singular (not invertible).
		\item Swapping two rows multiplies the determinant by $-1$.
		\item Multiplying a row by a scalar $k$ multiplies the determinant by $k$.
		\item Adding a multiple of one row to another row does not change the determinant.
		\item $\det(AB) = \det(A)\det(B)$ for any $n\times n$ matrices $A,B$.
		\item $\det(A^{T}) = \det(A)$.
	\end{itemize}
	
	\section*{\underline{Examples}$\downarrow$}
	
	\subsection*{Example 1: $2\times 2$ determinant}
	
	Compute the determinant of
	\[
	A =
	\begin{bmatrix}
		2 & -3 \\
		4 & 1
	\end{bmatrix}.
	\]
	
	Using the formula $\det(A) = ad - bc$, we have
	\[
	\det(A) = (2)(1) - (-3)(4)
	\]
	\[
	\det(A) = 2 - (-12)
	\]
	\[
	\det(A) = 2 + 12 = 14.
	\]
	
	\subsection*{Example 2: $3\times 3$ determinant (cofactor expansion)}
	
	Compute the determinant of
	\[
	B =
	\begin{bmatrix}
		1 & 2 & 3 \\
		0 & -1 & 4 \\
		2 & 1 & 0
	\end{bmatrix}.
	\]
	
	Expand along the first row:
	\[
	\det(B) =
	1
	\begin{vmatrix}
		-1 & 4 \\
		1 & 0
	\end{vmatrix}
	-
	2
	\begin{vmatrix}
		0 & 4 \\
		2 & 0
	\end{vmatrix}
	+
	3
	\begin{vmatrix}
		0 & -1 \\
		2 & 1
	\end{vmatrix}.
	\]
	
	Now evaluate each $2\times 2$ determinant:
	\[
	\begin{vmatrix}
		-1 & 4 \\
		1 & 0
	\end{vmatrix}
	= (-1)(0) - (4)(1) = -4,
	\]
	\[
	\begin{vmatrix}
		0 & 4 \\
		2 & 0
	\end{vmatrix}
	= (0)(0) - (4)(2) = -8,
	\]
	\[
	\begin{vmatrix}
		0 & -1 \\
		2 & 1
	\end{vmatrix}
	= (0)(1) - (-1)(2) = 2.
	\]
	
	Substitute back:
	\[
	\det(B) = 1(-4) - 2(-8) + 3(2)
	\]
	\[
	\det(B) = -4 + 16 + 6
	\]
	\[
	\det(B) = 18.
	\]
	
	\subsection*{Example 3: $3\times 3$ determinant (Sarrus' rule)}
	
	Compute the determinant of
	\[
	C =
	\begin{bmatrix}
		2 & 0 & 1 \\
		-1 & 3 & 2 \\
		4 & -2 & 1
	\end{bmatrix}.
	\]
	
	Using Sarrus' rule, write the first two columns again:
	\[
	\begin{array}{ccc|cc}
		2 & 0 & 1 & 2 & 0 \\
		-1 & 3 & 2 & -1 & 3 \\
		4 & -2 & 1 & 4 & -2
	\end{array}
	\]
	
	Sum of downward diagonals:
	\[
	(2)(3)(1) + (0)(2)(4) + (1)(-1)(-2) = 6 + 0 + 2 = 8.
	\]
	
	Sum of upward diagonals:
	\[
	(1)(3)(4) + (2)(2)(2) + (0)(-1)(1) = 12 + 8 + 0 = 20.
	\]
	
	Therefore,
	\[
	\det(C) = 8 - 20 = -12.
	\]
	
	\subsection*{Example 4: $4\times 4$ determinant with zeros}
	
	Compute the determinant of
	\[
	D =
	\begin{bmatrix}
		1 & 0 & 2 & -1 \\
		3 & 1 & 0 & 2 \\
		0 & -2 & 1 & 0 \\
		2 & 0 & 1 & 1
	\end{bmatrix}.
	\]
	
	Expand along the first row (because it has a zero):
	\[
	\det(D) =
	1 \cdot C_{11}
	+
	0 \cdot C_{12}
	+
	2 \cdot C_{13}
	+
	(-1) \cdot C_{14}.
	\]
	
	So
	\[
	\det(D) = 1 \cdot
	\begin{vmatrix}
		1 & 0 & 2 \\
		-2 & 1 & 0 \\
		0 & 1 & 1
	\end{vmatrix}
	+
	2 \cdot
	\begin{vmatrix}
		3 & 1 & 2 \\
		0 & -2 & 0 \\
		2 & 0 & 1
	\end{vmatrix}
	-
	1 \cdot
	\begin{vmatrix}
		3 & 1 & 0 \\
		0 & -2 & 1 \\
		2 & 0 & 1
	\end{vmatrix}.
	\]
	
	Call these three $3\times 3$ determinants $D_1$, $D_2$, and $D_3$.
	
	\newpage
	
	\paragraph{Compute $D_1$:}
	\[
	D_1 =
	\begin{vmatrix}
		1 & 0 & 2 \\
		-2 & 1 & 0 \\
		0 & 1 & 1
	\end{vmatrix}.
	\]
	
	Expand along the first row:
	\[
	D_1 =
	1
	\begin{vmatrix}
		1 & 0 \\
		1 & 1
	\end{vmatrix}
	-
	0
	\begin{vmatrix}
		-2 & 0 \\
		0 & 1
	\end{vmatrix}
	+
	2
	\begin{vmatrix}
		-2 & 1 \\
		0 & 1
	\end{vmatrix}.
	\]
	
	Compute the $2\times 2$ determinants:
	\[
	\begin{vmatrix}
		1 & 0 \\
		1 & 1
	\end{vmatrix}
	= (1)(1) - (0)(1) = 1,
	\]
	\[
	\begin{vmatrix}
		-2 & 1 \\
		0 & 1
	\end{vmatrix}
	= (-2)(1) - (1)(0) = -2.
	\]
	
	Thus
	\[
	D_1 = 1(1) + 2(-2) = 1 - 4 = -3.
	\]
	
	\newpage
	
	\paragraph{Compute $D_2$:}
	\[
	D_2 =
	\begin{vmatrix}
		3 & 1 & 2 \\
		0 & -2 & 0 \\
		2 & 0 & 1
	\end{vmatrix}.
	\]
	
	Expand along the second row (it has a zero):
	\[
	D_2 =
	0 \cdot
	\begin{vmatrix}
		1 & 2 \\
		0 & 1
	\end{vmatrix}
	-
	(-2) \cdot
	\begin{vmatrix}
		3 & 2 \\
		2 & 1
	\end{vmatrix}
	+
	0 \cdot
	\begin{vmatrix}
		3 & 1 \\
		2 & 0
	\end{vmatrix}.
	\]
	
	Only the middle term remains:
	\[
	D_2 = 2
	\begin{vmatrix}
		3 & 2 \\
		2 & 1
	\end{vmatrix}.
	\]
	
	Compute the $2\times 2$ determinant:
	\[
	\begin{vmatrix}
		3 & 2 \\
		2 & 1
	\end{vmatrix}
	= (3)(1) - (2)(2) = 3 - 4 = -1.
	\]
	
	So
	\[
	D_2 = 2(-1) = -2.
	\]
	
	\paragraph{Compute $D_3$:}
	\[
	D_3 =
	\begin{vmatrix}
		3 & 1 & 0 \\
		0 & -2 & 1 \\
		2 & 0 & 1
	\end{vmatrix}.
	\]
	
	Expand along the first row (because of the zero):
	\[
	D_3 =
	3
	\begin{vmatrix}
		-2 & 1 \\
		0 & 1
	\end{vmatrix}
	-
	1
	\begin{vmatrix}
		0 & 1 \\
		2 & 1
	\end{vmatrix}
	+
	0 \cdot
	\begin{vmatrix}
		0 & -2 \\
		2 & 0
	\end{vmatrix}.
	\]
	
	Compute the $2\times 2$ determinants:
	\[
	\begin{vmatrix}
		-2 & 1 \\
		0 & 1
	\end{vmatrix}
	= (-2)(1) - (1)(0) = -2,
	\]
	\[
	\begin{vmatrix}
		0 & 1 \\
		2 & 1
	\end{vmatrix}
	= (0)(1) - (1)(2) = -2.
	\]
	
	Thus
	\[
	D_3 = 3(-2) - 1(-2) = -6 + 2 = -4.
	\]
	
	\paragraph{Combine to get $\det(D)$:}
	\[
	\det(D) = 1\cdot D_1 + 2\cdot D_2 - 1\cdot D_3
	\]
	\[
	\det(D) = (-3) + 2(-2) - (-4)
	\]
	\[
	\det(D) = -3 - 4 + 4 = -3.
	\]
	
	So
	\[
	\det(D) = -3.
	\]
	
\end{document}